\section{Constant budgets and restricted partitions}

For $d\ge1$, define
\[
 p_d(n)=\#\{(a_1,\ldots,a_d)\in\N^d:
                         a_1+2a_2+\cdots+da_d=n\}.
\]
The independent choices of multiplicities give
\begin{equation}\label{eq:partition-series}
 \sum_{n\ge0}p_d(n)z^n=\prod_{i=1}^d(1-z^i)^{-1}.
\end{equation}
The relation to ordinals preserves every part, including the finite tail.

\begin{lemma}[A fixed-budget count]\label{lem:fixed}
Fix a nonnegative retained mass $c$ and take $S\ge c+d+1$.
Start the suffix at $\omega^d$ with any finite tail of total norm
(including $c$) at most $S$. At every replacement spend the constant
suffix budget $S-c$ and stop when the suffix is finite. The number of
replacements is $p_d(S-c)$. When $c=\norm(\pi)$ for a retained
CNF prefix $\pi$ of least exponent at least $d$, this is precisely
the corresponding segment starting at $\pi+\omega^d$ plus that tail.
\end{lemma}
\begin{proof}
The ordinals of suffix norm exactly $S-c$ below $\omega^d$ are
precisely
\[
 \omega^{d-1}a_d+\cdots+\omega a_2+a_1,\qquad
             \sum_{i=1}^d i a_i=S-c.
 \tag{*}
\]
Their number is $p_d(S-c)$. The first predecessor is the largest
such ordinal. In any later state of (*), the nonfinite coefficients
determine $a_1$ uniquely. Thus no second state of (*) has the same
nonfinite part. Removing the finite tail and taking the largest
predecessor of suffix norm $S-c$ selects exactly the next state of
(*) in decreasing ordinal order. All states are visited, ending at
the finite suffix $S-c$. There is one replacement into each state
and none after the last, proving the count. For $d=1$ there is one
state and one replacement.
\end{proof}

\begin{lemma}[A sufficient partition period]\label{lem:period}
For every $d,m\ge1$, put $Q_d=d!m^{d-1}$. Then
\[
             p_d(n+Q_d)\equiv p_d(n)\pmod m\quad(n\ge0).
\]
\end{lemma}
\begin{proof}
For $d=1$, $p_1(n)=1$. For $d>1$, splitting by the number
of parts of size $d$ gives
$p_d(n)=\sum_{j=0}^{\lfloor n/d\rfloor}p_{d-1}(n-dj)$.
Here $Q_d/d=mQ_{d-1}$. In the sum for $n+Q_d$, shifting the
index by $mQ_{d-1}$ recovers the old sum. The additional
$mQ_{d-1}$ summands form $m$ groups of $Q_{d-1}$ consecutive
indices in the periodic sequence
$j\mapsto p_{d-1}(n+Q_d-dj)\pmod m$.
Its period divides $Q_{d-1}$ by induction, so the extra sum is zero
modulo $m$.
\end{proof}

The period is deliberately a sufficient bound, not an optimal one.
For $d=2$ the count is $\lfloor n/2\rfloor+1$, and $2m$
is the period supplied by the lemma.

\section{A transfer theorem for congruent growing budgets}

For the numerical statement, a retained mass $c\ge0$ means that a
suffix of norm $n$ has total norm $c+n$, and a total budget $b$ leaves
suffix budget $b-c$. It need not be the norm of an ordinal prefix.
If it is realized by an actual retained CNF prefix, the numerical
process agrees with that source segment. A budget schedule may increase
with the current total norm; it must always afford that norm.

\begin{theorem}[Congruent-budget transfer]\label{thm:transfer}
Consider deletion of one suffix block $\omega^d$, with retained
numerical mass $c\ge0$ and initial total norm at least $c+d+1$,
under budgets at least the current total norm. Suppose
every budget used in this segment is congruent to one integer $S$
modulo $Q_d=d!m^{d-1}$, where $S\ge c$. Then its complete segment
count is congruent to $p_d(S-c)$ modulo $m$.
\end{theorem}
\begin{proof}
Induct on $d$. A degree-one block takes one replacement, as does
$p_1(S-c)$. For the induction step let $b$ be the first actual
total budget. The suffix of its greatest predecessor is the greedy
Cantor polynomial of norm $b-c$ and degree less than $d$.
Write its coefficients as $a_{d-1},\ldots,a_0$, with the final
finite coefficient included.

Thereafter each of its nonfinite monomial units is deleted by a
lower-degree block. For a degree-$i$ unit the retained prefix norm
has the form
\[
 c_i=c+\sum_{j>i}(j+1)a_j+(i+1)a,\qquad 0\le a<a_i.
 \tag{**}
\]
The order of those units is fixed by the Cantor order; growing
budgets affect their lower suffixes and durations, not the number
of these fixed higher monomial units. Since $Q_i\mid Q_d$, every
budget in a child segment has the residue of $b$ modulo $Q_i$.
Induction makes its count congruent to $p_i(b-c_i)$.
These arguments are nonnegative: each $c_i$ consists of $c$ plus
the norm of a prefix of the first suffix polynomial and is at most $b$.

Consequently the whole count, including its initial replacement,
has the same residue as the constant-$b$ process with exactly that
initial polynomial and those prefixes (**). By
Lemma~\ref{lem:fixed} this is $p_d(b-c)$.
Since $b\equiv S\pmod {Q_d}$, Lemma~\ref{lem:period} gives
$p_d(b-c)\equiv p_d(S-c)\pmod m$. If $b-c<S-c$,
apply the period lemma in the reverse direction by adding the
appropriate nonnegative multiple; no negative partition index is used.
\end{proof}

The proof covers the whole count, not just its leading contribution.
It does not replace the increasing operational budgets by a constant
control: the latter is only a device for computing one residue.

\begin{corollary}[A fixed polynomial block under the standard control]
\label{cor:fixed-polynomial-residue}
Fix $c,d,m$ and put $Q_d=d!m^{d-1}$. Under the control $G$,
the complete count $E_c^d(t)$ of one $\omega^d$ block with retained
mass $c$ and total input norm $t$ satisfies
\[
 E_c^d(t)\equiv
 p_d\bigl((r_{Q_d}-c)\bmod Q_d\bigr)\pmod m
 \quad\text{if }t\ge\max(c+d+1,Q_d).
\]
The finitely many smaller valid inputs can be evaluated exactly.
\end{corollary}
\begin{proof}
All later norms are at least $t$, so every budget has residue
$r_{Q_d}$ by Theorem~\ref{thm:envelope}. Apply the transfer theorem
with $S=c+((r_{Q_d}-c)\bmod Q_d)\ge c$.
Exact evaluation handles the remaining finite set.
\end{proof}

